\section{References}

\subsection*{Sacred Scripture}

\begin{itemize}
\item \textbf{The Holy Bible: Douay--Rheims Version (Challoner Revision).} Electronic text of Bishop Richard Challoner's 1749--1752 revision, Project Gutenberg eBook 1581 (updated 2023; exemplar printing not individually identified; public domain). Cited at Gen 12:1--3; Deut 7:6--8, 10:14--19; Ps 8:4--6; Wis 11:21, 25, 27; Isa 2:2--3, 49:6; Amos 9:7; Matt 1:1, 16--17, 13:55--57, 16:15, 25:35--36, 40; Mark 6:3--5; Luke 2:1--2, 3:1--2, 23, 38; John 1:14, 45--46, 4:24, 14:6, 20:27--28, 21:9, 12--13, 15--17; Acts 4:12, 17:27--28; Rom 9:4--5, 11:17--18, 28--29, 33; 1 Cor 1:23; Gal 4:4--5; 1 Tim 2:4--5; Heb 1:1--2, 2:16; 1 John 4:19--21; Apoc 2:17, 5:9.
\end{itemize}

\subsection*{Patristic and Scholastic Primary Witnesses}

\begin{itemize}
\item \textbf{Athanasius of Alexandria, St.} \textit{On the Incarnation of the Word} (\latin{De Incarnatione Verbi Dei}). Translated by Archibald Robertson. In \textit{A Select Library of Nicene and Post-Nicene Fathers of the Christian Church}, Second Series, vol.\ 4, edited by Philip Schaff, Henry Wace, and Archibald Robertson. New York: The Christian Literature Company; Oxford and London: Parker \& Company, 1892. (Public domain; cited at \S\S\,8.1, 9.3, 41.5, 42.7, 43.1, 3--4, 44.2, 5--8, 54.2--3.)
\item \textbf{Augustine of Hippo, St.} \textit{The City of God} (\latin{De Civitate Dei}). Translated by Marcus Dods (with George Wilson and J.\ J.\ Smith; Book X translated by Marcus Dods). Vol.\ 1. Edinburgh: T.\ \& T.\ Clark, 1871. (Public domain; Project Gutenberg eBook 45304; cited at X.5--6, 28--29, 32.)
\item \textbf{Augustine of Hippo, St.} \textit{Tractates on the Gospel of John} (\latin{In Iohannis Evangelium Tractatus}). Translated by John Gibb and James Innes. In \textit{Nicene and Post-Nicene Fathers}, First Series, vol.\ 7, edited by Philip Schaff. Christian Classics Ethereal Library dated electronic text (2026-09-17; file header cites New York: Christian Literature Publishing Co., 1886; identity with a specific historical printing not asserted; public domain; cited at Tractate 80, \S\,3, and Tractate 123, \S\S\,2, 5.)
\item \textbf{Gregory of Nazianzus, St.} \textit{Oration 38: On the Theophany, or Birthday of Christ}. Translated by Charles Gordon Browne and James Edward Swallow. In \textit{Nicene and Post-Nicene Fathers}, Second Series, vol.\ 7, edited by Philip Schaff and Henry Wace. Christian Classics Ethereal Library dated electronic text (2026-10-06; file header cites New York: Christian Literature Publishing Co., 1893; not collated as a facsimile; public domain; cited at \S\S\,11, 13--14.)
\item \textbf{Irenaeus of Lyons, St.} \textit{Against Heresies} (\latin{Adversus Haereses}). Translated by Alexander Roberts and William Rambaut. In \textit{The Ante-Nicene Fathers}, vol.\ 1, edited by Alexander Roberts, James Donaldson, and A.\ Cleveland Coxe. Christian Classics Ethereal Library dated electronic text (2026-10-13; exact reprint date not asserted) for Books III--IV, supplemented by the checked Buffalo: Christian Literature Publishing Co., 1887 printing for Book V, preface. (Public domain; cited at III.18.1, 7; III.22.1--2, 4; IV.22.2; V, preface.)
\item \textbf{John Chrysostom, St.} \textit{Homilies on the Gospel of St.\ Matthew}. Translated by George Prevost and M.\ B.\ Riddle. In \textit{Nicene and Post-Nicene Fathers}, First Series, vol.\ 10, edited by Philip Schaff. Christian Classics Ethereal Library dated electronic text (2026-09-17; file header cites New York: Christian Literature Publishing Co., 1886; identity with a specific historical printing not asserted; public domain; cited at Homily 50, \S\S\,4--5.)
\item \textbf{Leo the Great, St.} \textit{Letters} and \textit{Sermons}. Translated by Charles Lett Feltoe. In \textit{Nicene and Post-Nicene Fathers}, Second Series, vol.\ 12, edited by Philip Schaff and Henry Wace. Christian Classics Ethereal Library dated electronic text (2026-09-28; exact print impression not asserted) for \textit{Letter 28} and \textit{Sermon 74}, supplemented by checked excerpts of the New York, 1895 printing for \textit{Sermons 21} and \textit{54}. (Public domain; cited at \textit{Letter 28} [\textit{Tome to Flavian}], \S\S\,3--4; \textit{Sermon 21}, \S\S\,1--2; \textit{Sermon 54}, \S\,4; \textit{Sermon 74}, \S\,2.)
\item \textbf{Thomas Aquinas, St.} \textit{De ente et essentia}. Latin text edited by Ludwig Baur and emended by Josef Koch. M\"unster: Aschendorff, 1933. (Public domain; cited at cc.\,2--3, 5.)
\item \textbf{Thomas Aquinas, St.} \textit{Summa Theologiae}, Prima Pars and Tertia Pars. Translated by the Fathers of the English Dominican Province. Electronic text of the Benziger Brothers translation, Project Gutenberg eBooks 17611 (2006, rev.\ 2021) and 19950 (2006; exact print-edition year not inferred from electronic releases; public domain). Cited at I, q.\,3, a.\,5; q.\,8, a.\,1; q.\,14, a.\,11; q.\,22, aa.\,2--4; q.\,29, aa.\,1, 3; q.\,44, a.\,1; q.\,85, a.\,1; q.\,86, a.\,1; q.\,104, a.\,1; III, q.\,1, a.\,1 (including Volusianus, \textit{Ep.}\ 135, and Augustine, \textit{Ep.}\ 137); q.\,2, aa.\,1--2 (including the Council of Chalcedon); q.\,4, aa.\,4--5 (including St.\ John Damascene, \textit{De Fide Orthodoxa} III.11); q.\,60, a.\,4; q.\,61, a.\,1; q.\,64, a.\,7; q.\,75, a.\,1; q.\,76, a.\,3.
\end{itemize}

\subsection*{Conciliar and Magisterial Documents (Doctrinal Paraphrase Only)}

\begin{itemize}
\item \textbf{Catholic Church.} \textit{Catechismus Catholicae Ecclesiae} (\textit{Catechism of the Catholic Church}), \latin{editio typica}. Vatican City: Libreria Editrice Vaticana, 1997. (Cataloged rights-restricted witness; cited in independent doctrinal paraphrase at \S\S\,1084--1085, 1145--1152.)
\item \textbf{Congregation for the Doctrine of the Faith.} Declaration \textit{Dominus Iesus: On the Unicity and Salvific Universality of Jesus Christ and the Church} (6 August 2000). \latin{Acta Apostolicae Sedis} 92 (2000): 742--765. (Cataloged rights-restricted witness; cited in independent doctrinal paraphrase at \S\S\,6, 13--15.)
\item \textbf{John Paul II, Pope St.} Encyclical Letter \textit{Redemptoris missio: On the Permanent Validity of the Church's Missionary Mandate} (7 December 1990). \latin{Acta Apostolicae Sedis} 83 (1991): 249--340. (Cataloged rights-restricted witness; cited in independent doctrinal paraphrase at \S\S\,5--6, 28.)
\item \textbf{John Paul II, Pope St.} Encyclical Letter \textit{Fides et ratio: On the Relationship between Faith and Reason} (14 September 1998). \latin{Acta Apostolicae Sedis} 91 (1999): 5--88. (Cataloged rights-restricted witness; cited in independent doctrinal paraphrase at \S\S\,7, 11--12.)
\item \textbf{Second Vatican Council.} Dogmatic Constitution on the Church \textit{Lumen gentium} (21 November 1964). \latin{Acta Apostolicae Sedis} 57 (1965): 5--71. (Cataloged rights-restricted witness; cited in independent doctrinal paraphrase at \S\,16.)
\item \textbf{Second Vatican Council.} Declaration on the Relation of the Church to Non-Christian Religions \textit{Nostra aetate} (28 October 1965). \latin{Acta Apostolicae Sedis} 58 (1966): 740--744. (Cataloged rights-restricted witness; cited in independent doctrinal paraphrase at \S\S\,2, 4.)
\item \textbf{Second Vatican Council.} Dogmatic Constitution on Divine Revelation \textit{Dei Verbum} (18 November 1965). \latin{Acta Apostolicae Sedis} 58 (1966): 817--835. (Cataloged rights-restricted witness; cited in independent doctrinal paraphrase at \S\S\,2--4, 14--16.)
\item \textbf{Second Vatican Council.} Pastoral Constitution on the Church in the Modern World \textit{Gaudium et spes} (7 December 1965). \latin{Acta Apostolicae Sedis} 58 (1966): 1025--1120. (Cataloged rights-restricted witness; cited in independent doctrinal paraphrase at \S\,22.)
\end{itemize}

\subsection*{Historical and Modern Interlocutors (Conceptual Engagement)}

\begin{itemize}
\item \textbf{Boethius, Anicius Manlius Severinus.} \textit{Contra Eutychen et Nestorium} (\textit{De duabus naturis}). Referenced via St.\ Thomas Aquinas, \textit{ST} I, q.\,29, a.\,1, for the classical definition of person.
\item \textbf{Chalmers, Thomas.} \textit{A Series of Discourses on the Christian Revelation, Viewed in Connection with the Modern Astronomy} (1817). Engaged conceptually as a classic theological response to the astronomical objection.
\item \textbf{Celsus (via Origen, \textit{Contra Celsum}).} Engaged conceptually for the second-century philosophical objection to divine descent in one corner of the earth.
\item \textbf{Duns Scotus, Bl.\ John, and Francisco Su\'arez.} Referenced historically for alternative Catholic scholastic accounts of individuation (\latin{haecceitas} and whole-entity individuation) contrasted with St.\ Thomas's designated matter.
\item \textbf{Hick, John.} \textit{An Interpretation of Religion: Human Responses to the Transcendent} (1989). Engaged conceptually for the modern pluralist hypothesis regarding religious diversity and the geography of birth.
\item \textbf{Kant, Immanuel.} \textit{Religion within the Boundaries of Mere Reason} (1793). Engaged conceptually for treating historical ecclesiastical faith as a provisional vehicle for pure moral religion.
\item \textbf{Kierkegaard, S\o ren.} \textit{Philosophical Fragments} (1844). Engaged conceptually for engaging Lessing's historical ditch through the absolute paradox of the Eternal in time.
\item \textbf{Lessing, Gotthold Ephraim.} \textit{On the Proof of the Spirit and of Power} (1777). Engaged conceptually for the epistemological ditch between contingent historical truths and necessary truths of reason.
\item \textbf{Paine, Thomas.} \textit{The Age of Reason} (1794). Engaged conceptually for the Enlightenment astronomical objection against a singular earthly Incarnation.
\item \textbf{Spinoza, Baruch.} \textit{Theologico-Political Treatise} (1670). Engaged conceptually for the rationalist dissolution of biblical election into natural political fortune.
\end{itemize}
